% !TEX root = userguide.tex

\def\headdate{12th November 2021}

\section{Blended meshes}
\label{sec:blendexamples}

In Sec.~\ref{sec:blendedmeshes} blended meshes are introduced to
easily couple variables with different orders of interpolation. 
In this section some examples using blended meshes are given.

In the following table an overview of the blended mesh examples in \tfem\ are given\medskip

\begin{tabular}{lll}
Example & Section & remarks \\\hline
 \texttt{sphere31.f90} & \ref{sec:stokes31} & Taylor-Hood  $Q_2/Q_1$  \\
 \texttt{sphere32.f90} & \ref{sec:stokes31} & Taylor-Hood  $Q_4/Q_3$  \\
 \texttt{mixed\_stokes1.f90} & \ref{sec:mixedcontraction} & mixed Stokes Taylor-Hood $Q_2/Q_1$ with $Q_p$ stress \\
 \texttt{mixed\_stokes2.f90} & \ref{sec:mixedcontraction} & mixed Stokes Taylor-Hood $Q_2/Q_1$ with $Q_1\text{-iso-}Q_p$ stress \\
\end{tabular}


\subsection{Stokes flow around a sphere in a cylindrical tube}
\label{sec:stokes31}

The examples in this section are similar to the \texttt{sphere1.f90} problem 
of Sec.~\ref{sec:sphere}.
The first example \texttt{sphere31.f90} uses the same  $Q_2/Q_1$  Taylor-Hood
interpolation for velocity/pressure as \texttt{sphere1.f90}. However,
the velocity and pressure are now defined on their own mesh using a blended
mesh. The second example \texttt{sphere32.f90} uses $Q_4/Q_3$ Taylor-Hood
interpolation on a blended mesh. The latter example would be difficult to
implement using a single shape element because of the different number and
positions of the edge nodes for velocity and pressure.
